\expandafter\gdef\csname SNside1\endcsname{\node[anchor=north west,inner sep=0pt,text width=48mm,font=\SideFont,align=left] (SNsidebox) at ([xshift=138mm,yshift=-407.15pt]current page.north west) {\textcolor{SNink}{\textbf{课程连接 · 两条路径}}\par\vspace{1.5mm}目录路径解决名称定位，文件映射解决逻辑内容到磁盘块的定位，两者的I/O可能分别发生。};\pgfextracty{\SNsideheight}{\pgfpointdiff{\pgfpointanchor{SNsidebox}{south}}{\pgfpointanchor{SNsidebox}{north}}}\typeout{SNHEIGHT:1:19:\the\SNsideheight}}
\expandafter\gdef\csname SNside2\endcsname{\node[anchor=north west,inner sep=0pt,text width=48mm,font=\SideFont,align=left] (SNsidebox) at ([xshift=138mm,yshift=-95.00pt]current page.north west) {\textcolor{SNink}{\textbf{错项排除 · 同号不一定同文件}}\par\vspace{1.5mm}inode号要在所属文件系统内解释。两个不同分区上都出现inode10，不足以证明它们指向同一文件；同一文件系统中，两个名字解析到同一个inode才表示这种对象共享。判题时先确定命名空间，再比较数字。};\pgfextracty{\SNsideheight}{\pgfpointdiff{\pgfpointanchor{SNsidebox}{south}}{\pgfpointanchor{SNsidebox}{north}}}\typeout{SNHEIGHT:2:1:\the\SNsideheight}}
\expandafter\gdef\csname SNside3\endcsname{\node[anchor=north west,inner sep=0pt,text width=48mm,font=\SideFont,align=left] (SNsidebox) at ([xshift=138mm,yshift=-95.00pt]current page.north west) {\textcolor{SNink}{\textbf{最小反例 · 关闭只撤销本次使用引用}}\par\vspace{1.5mm}一个文件有名字a，进程打开后再关闭，目录里的a仍存在，文件数据也应保留。close减少的是打开使用关系，unlink取消的是名字关系。只有满足没有硬链接、没有有效使用引用等条件，才讨论对象最终回收。};\pgfextracty{\SNsideheight}{\pgfpointdiff{\pgfpointanchor{SNsidebox}{south}}{\pgfpointanchor{SNsidebox}{north}}}\typeout{SNHEIGHT:3:2:\the\SNsideheight}}
\expandafter\gdef\csname SNside4\endcsname{\node[anchor=north west,inner sep=0pt,text width=48mm,font=\SideFont,align=left] (SNsidebox) at ([xshift=138mm,yshift=-95.00pt]current page.north west) {\textcolor{SNink}{\textbf{辨析 · 同一文件，多次打开}}\par\vspace{1.5mm}图按课件模型把指针记在进程表、文件打开计数记在系统表。具体系统的表项划分与偏移共享规则需按题目说明判断。};\pgfextracty{\SNsideheight}{\pgfpointdiff{\pgfpointanchor{SNsidebox}{south}}{\pgfpointanchor{SNsidebox}{north}}}\typeout{SNHEIGHT:4:24:\the\SNsideheight}}
\expandafter\gdef\csname SNside5\endcsname{\node[anchor=north west,inner sep=0pt,text width=48mm,font=\SideFont,align=left] (SNsidebox) at ([xshift=138mm,yshift=-306.85pt]current page.north west) {\textcolor{SNink}{\textbf{错项排除 · 权限集合不能压成一个编号}}\par\vspace{1.5mm}每类用户可能同时有读、写等多种权限；若有5个独立许可就要表达其全部子集，需要5位。给用户类别和某一操作编码，只能描述一次选择，不能保存每一类用户各自的完整许可组合，二者信息量不同。};\pgfextracty{\SNsideheight}{\pgfpointdiff{\pgfpointanchor{SNsidebox}{south}}{\pgfpointanchor{SNsidebox}{north}}}\typeout{SNHEIGHT:5:4:\the\SNsideheight}}
\expandafter\gdef\csname SNside6\endcsname{\node[anchor=north west,inner sep=0pt,text width=48mm,font=\SideFont,align=left] (SNsidebox) at ([xshift=138mm,yshift=-95.00pt]current page.north west) {\textcolor{SNink}{\textbf{辨析 · 目录权限不等于内容权限}}\par\vspace{1.5mm}能遍历目录与能读文件内容属于不同操作。修改文件和修改目录项也不同，必须按被访问对象及具体动作验证。};\pgfextracty{\SNsideheight}{\pgfpointdiff{\pgfpointanchor{SNsidebox}{south}}{\pgfpointanchor{SNsidebox}{north}}}\typeout{SNHEIGHT:6:18:\the\SNsideheight}}
\expandafter\gdef\csname SNside7\endcsname{\node[anchor=north west,inner sep=0pt,text width=48mm,font=\SideFont,align=left] (SNsidebox) at ([xshift=138mm,yshift=-95.00pt]current page.north west) {\textcolor{SNink}{\textbf{课程连接 · 两个索引属于不同层次}}\par\vspace{1.5mm}逻辑记录索引把关键字对应到记录或文件内偏移；物理索引把逻辑块号对应到磁盘块。数据结构中的查找复杂度可解释前一层，操作系统的间接索引层数解释后一层。遇到``索引''先写它的输入与输出对象。};\pgfextracty{\SNsideheight}{\pgfpointdiff{\pgfpointanchor{SNsidebox}{south}}{\pgfpointanchor{SNsidebox}{north}}}\typeout{SNHEIGHT:7:5:\the\SNsideheight}}
\expandafter\gdef\csname SNside8\endcsname{\node[anchor=north west,inner sep=0pt,text width=48mm,font=\SideFont,align=left] (SNsidebox) at ([xshift=138mm,yshift=-95.00pt]current page.north west) {\textcolor{SNink}{\textbf{辨析 · 两次映射先后不同}}\par\vspace{1.5mm}先由记录组织确定文件内逻辑位置，再由文件系统映射逻辑块到物理块；物理连续不会自动让链式记录支持按位序直接定位。};\pgfextracty{\SNsideheight}{\pgfpointdiff{\pgfpointanchor{SNsidebox}{south}}{\pgfpointanchor{SNsidebox}{north}}}\typeout{SNHEIGHT:8:26:\the\SNsideheight}}
\expandafter\gdef\csname SNside9\endcsname{\node[anchor=north west,inner sep=0pt,text width=48mm,font=\SideFont,align=left] (SNsidebox) at ([xshift=138mm,yshift=-354.47pt]current page.north west) {\textcolor{SNink}{\textbf{辨析 · 两次映射先后不同}}\par\vspace{1.5mm}先由记录组织确定文件内逻辑位置，再由文件系统映射逻辑块到物理块；物理连续不会自动让链式记录支持按位序直接定位。};\pgfextracty{\SNsideheight}{\pgfpointdiff{\pgfpointanchor{SNsidebox}{south}}{\pgfpointanchor{SNsidebox}{north}}}\typeout{SNHEIGHT:9:26:\the\SNsideheight}}
\expandafter\gdef\csname SNside10\endcsname{\node[anchor=north west,inner sep=0pt,text width=48mm,font=\SideFont,align=left] (SNsidebox) at ([xshift=138mm,yshift=-95.00pt]current page.north west) {\textcolor{SNink}{\textbf{错项排除 · FAT少读数据块仍需追链}}\par\vspace{1.5mm}FAT驻留内存时，定位后面簇不必逐个读取前面数据块，但仍要沿内存中的链获取后继号。其优势是减少磁盘I/O，不能直接推成按逻辑下标常数时间定位；表规模还由整个受管卷的簇数决定。};\pgfextracty{\SNsideheight}{\pgfpointdiff{\pgfpointanchor{SNsidebox}{south}}{\pgfpointanchor{SNsidebox}{north}}}\typeout{SNHEIGHT:10:6:\the\SNsideheight}}
\expandafter\gdef\csname SNside11\endcsname{\node[anchor=north west,inner sep=0pt,text width=48mm,font=\SideFont,align=left] (SNsidebox) at ([xshift=138mm,yshift=-95.00pt]current page.north west) {\textcolor{SNink}{\textbf{不变量 · 容量与索引开销分开累计}}\par\vspace{1.5mm}一个索引块可装B除以指针字节数个地址。计算最大文件长度时只累加可指向的数据块，索引块本身另占外存。若题目转问总空间消耗，才把实际使用的各级索引块加回，两个问题不能共用一个总数。};\pgfextracty{\SNsideheight}{\pgfpointdiff{\pgfpointanchor{SNsidebox}{south}}{\pgfpointanchor{SNsidebox}{north}}}\typeout{SNHEIGHT:11:7:\the\SNsideheight}}
\expandafter\gdef\csname SNside12\endcsname{\node[anchor=north west,inner sep=0pt,text width=48mm,font=\SideFont,align=left] (SNsidebox) at ([xshift=138mm,yshift=-95.00pt]current page.north west) {\textcolor{SNink}{\textbf{不变量 · 缓存只删去对应访问}}\par\vspace{1.5mm}目录内容已在内存，只免去读取该目录内容；若目标inode和数据尚未缓存，仍要分别取到。写访盘路径时为每个节点标``已缓存/未缓存''，再计边上的实际读取，避免一个命中条件被扩张成整条路径全命中。};\pgfextracty{\SNsideheight}{\pgfpointdiff{\pgfpointanchor{SNsidebox}{south}}{\pgfpointanchor{SNsidebox}{north}}}\typeout{SNHEIGHT:12:8:\the\SNsideheight}}
\expandafter\gdef\csname SNside13\endcsname{\node[anchor=north west,inner sep=0pt,text width=48mm,font=\SideFont,align=left] (SNsidebox) at ([xshift=138mm,yshift=-95.00pt]current page.north west) {\textcolor{SNink}{\textbf{辨析 · 共享不是复制}}\par\vspace{1.5mm}多个目录项指向共同文件，更新可通过不同名字看到。删除一个目录项只取消一条关联，回收还需检查链接与打开使用状态。};\pgfextracty{\SNsideheight}{\pgfpointdiff{\pgfpointanchor{SNsidebox}{south}}{\pgfpointanchor{SNsidebox}{north}}}\typeout{SNHEIGHT:13:15:\the\SNsideheight}}
\expandafter\gdef\csname SNside14\endcsname{\node[anchor=north west,inner sep=0pt,text width=48mm,font=\SideFont,align=left] (SNsidebox) at ([xshift=138mm,yshift=-95.00pt]current page.north west) {\textcolor{SNink}{\textbf{最小反例 · 删一个名字不等于删文件}}\par\vspace{1.5mm}同一文件inode为i7，目录项a和b都指向它。删除a后仍有b这一硬链接，文件内容继续可达。把a叫作原文件不赋予它特殊地位；硬链接之间对称，不能说删除``原名字''就必然使另一个硬链接失效。\par\smallskip\begin{tikzpicture}[x=1mm,y=1mm,>=Stealth,node distance=5mm]
\node[draw=SNink!50,fill=SNblue,rounded corners=1mm,text width=34mm,align=center,inner sep=1.2mm,font=\fontsize{7.5}{9}\selectfont] (m0) at (0,0) {a指向i7};
\node[draw=SNink!50,fill=SNblue,rounded corners=1mm,text width=34mm,align=center,inner sep=1.2mm,font=\fontsize{7.5}{9}\selectfont] (m1) at (0,-11) {b指向i7};
\node[draw=SNink!50,fill=SNblue,rounded corners=1mm,text width=34mm,align=center,inner sep=1.2mm,font=\fontsize{7.5}{9}\selectfont] (m2) at (0,-22) {删a留b};
\draw[->,SNink] (m0) edge[bend left=40] (m2);
\draw[->,SNink] (m1) -- (m2);
\end{tikzpicture}
\par{\fontsize{7.5}{9}\selectfont 两个目录项对应同一inode；删除名字a后，名字b仍可引用文件。}};\pgfextracty{\SNsideheight}{\pgfpointdiff{\pgfpointanchor{SNsidebox}{south}}{\pgfpointanchor{SNsidebox}{north}}}\typeout{SNHEIGHT:14:0:\the\SNsideheight}}
\expandafter\gdef\csname SNside15\endcsname{\node[anchor=north west,inner sep=0pt,text width=48mm,font=\SideFont,align=left] (SNsidebox) at ([xshift=138mm,yshift=-420.41pt]current page.north west) {\textcolor{SNink}{\textbf{辨析 · 符号链接认路径}}\par\vspace{1.5mm}即使另一硬链接bbb仍可访问文件，删除aaa也会使保存C:/User1/aaa的ccc成为悬空链接。};\pgfextracty{\SNsideheight}{\pgfpointdiff{\pgfpointanchor{SNsidebox}{south}}{\pgfpointanchor{SNsidebox}{north}}}\typeout{SNHEIGHT:15:25:\the\SNsideheight}}
\expandafter\gdef\csname SNside16\endcsname{\node[anchor=north west,inner sep=0pt,text width=48mm,font=\SideFont,align=left] (SNsidebox) at ([xshift=138mm,yshift=-95.00pt]current page.north west) {\textcolor{SNink}{\textbf{辨析 · 按职责识别不同分层}}\par\vspace{1.5mm}名称查找、权限检查、记录地址、物理块地址和设备命令分别处理。教材与课件可合并或拆分层名，职责关系才是判断依据。};\pgfextracty{\SNsideheight}{\pgfpointdiff{\pgfpointanchor{SNsidebox}{south}}{\pgfpointanchor{SNsidebox}{north}}}\typeout{SNHEIGHT:16:17:\the\SNsideheight}}
\expandafter\gdef\csname SNside17\endcsname{\node[anchor=north west,inner sep=0pt,text width=48mm,font=\SideFont,align=left] (SNsidebox) at ([xshift=138mm,yshift=-95.00pt]current page.north west) {\textcolor{SNink}{\textbf{错项排除 · 位图长度取决于管理范围}}\par\vspace{1.5mm}管理一百万个块，即使当前只剩一个空闲块，位图仍须为每个块保留一位。空闲块数量影响位值分布，不缩短受管对象集合。位图内相对位置算完后再加其起始磁盘块号，不能先从受管块号减去起点。};\pgfextracty{\SNsideheight}{\pgfpointdiff{\pgfpointanchor{SNsidebox}{south}}{\pgfpointanchor{SNsidebox}{north}}}\typeout{SNHEIGHT:17:9:\the\SNsideheight}}
\expandafter\gdef\csname SNside18\endcsname{\node[anchor=north west,inner sep=0pt,text width=48mm,font=\SideFont,align=left] (SNsidebox) at ([xshift=138mm,yshift=-95.00pt]current page.north west) {\textcolor{SNink}{\textbf{辨析 · 承载块交付前先读目录}}\par\vspace{1.5mm}当前组最后的空闲块还保存下一组信息。必须先把组目录装入内存，再把该块用于文件数据，才能保住后续空闲块入口。};\pgfextracty{\SNsideheight}{\pgfpointdiff{\pgfpointanchor{SNsidebox}{south}}{\pgfpointanchor{SNsidebox}{north}}}\typeout{SNHEIGHT:18:16:\the\SNsideheight}}
\expandafter\gdef\csname SNside19\endcsname{\node[anchor=north west,inner sep=0pt,text width=48mm,font=\SideFont,align=left] (SNsidebox) at ([xshift=138mm,yshift=-95.00pt]current page.north west) {\textcolor{SNink}{\textbf{课程连接 · 统一接口可以隐藏多种布局}}\par\vspace{1.5mm}VFS把查找、打开、读写等操作抽象成共同接口，再由具体文件系统实现。它和面向接口编程相似：调用约定统一，并不要求内部数据结构一致。不能因为应用都调用read，就推断各磁盘格式具有相同的inode字段。};\pgfextracty{\SNsideheight}{\pgfpointdiff{\pgfpointanchor{SNsidebox}{south}}{\pgfpointanchor{SNsidebox}{north}}}\typeout{SNHEIGHT:19:10:\the\SNsideheight}}
\expandafter\gdef\csname SNside20\endcsname{\node[anchor=north west,inner sep=0pt,text width=48mm,font=\SideFont,align=left] (SNsidebox) at ([xshift=138mm,yshift=-95.00pt]current page.north west) {\textcolor{SNink}{\textbf{课程连接 · 统一接口可以隐藏多种布局}}\par\vspace{1.5mm}VFS把查找、打开、读写等操作抽象成共同接口，再由具体文件系统实现。它和面向接口编程相似：调用约定统一，并不要求内部数据结构一致。不能因为应用都调用read，就推断各磁盘格式具有相同的inode字段。};\pgfextracty{\SNsideheight}{\pgfpointdiff{\pgfpointanchor{SNsidebox}{south}}{\pgfpointanchor{SNsidebox}{north}}}\typeout{SNHEIGHT:20:10:\the\SNsideheight}}
\expandafter\gdef\csname SNside21\endcsname{\node[anchor=north west,inner sep=0pt,text width=48mm,font=\SideFont,align=left] (SNsidebox) at ([xshift=138mm,yshift=-95.00pt]current page.north west) {\textcolor{SNink}{\textbf{不变量 · 文件数量由最紧资源约束}}\par\vspace{1.5mm}每个图像既需要一个inode，也需要足够数据簇。应分别计算inode区能容纳的个数和数据区能容纳的个数，再取较小者；若还要间接索引簇或目录空间，应按题设另计，不能只用剩余字节除文件长度。};\pgfextracty{\SNsideheight}{\pgfpointdiff{\pgfpointanchor{SNsidebox}{south}}{\pgfpointanchor{SNsidebox}{north}}}\typeout{SNHEIGHT:21:11:\the\SNsideheight}}
\expandafter\gdef\csname SNside22\endcsname{\node[anchor=north west,inner sep=0pt,text width=48mm,font=\SideFont,align=left] (SNsidebox) at ([xshift=138mm,yshift=-95.00pt]current page.north west) {\textcolor{SNink}{\textbf{不变量 · 文件数量由最紧资源约束}}\par\vspace{1.5mm}每个图像既需要一个inode，也需要足够数据簇。应分别计算inode区能容纳的个数和数据区能容纳的个数，再取较小者；若还要间接索引簇或目录空间，应按题设另计，不能只用剩余字节除文件长度。};\pgfextracty{\SNsideheight}{\pgfpointdiff{\pgfpointanchor{SNsidebox}{south}}{\pgfpointanchor{SNsidebox}{north}}}\typeout{SNHEIGHT:22:11:\the\SNsideheight}}
\expandafter\gdef\csname SNside23\endcsname{\node[anchor=north west,inner sep=0pt,text width=48mm,font=\SideFont,align=left] (SNsidebox) at ([xshift=138mm,yshift=-95.00pt]current page.north west) {\textcolor{SNink}{\textbf{课程连接 · 检查顺序}}\par\vspace{1.5mm}区分目录查找、inode读取、索引块查询和数据块访问，再按题目缓存条件计数。};\pgfextracty{\SNsideheight}{\pgfpointdiff{\pgfpointanchor{SNsidebox}{south}}{\pgfpointanchor{SNsidebox}{north}}}\typeout{SNHEIGHT:23:21:\the\SNsideheight}}
\expandafter\gdef\csname SNside24\endcsname{\node[anchor=north west,inner sep=0pt,text width=48mm,font=\SideFont,align=left] (SNsidebox) at ([xshift=138mm,yshift=-95.00pt]current page.north west) {\textcolor{SNink}{\textbf{课程连接 · 回忆层次}}\par\vspace{1.5mm}从路径到目录项，再到索引节点和数据块；链接数与打开位置分别维护。};\pgfextracty{\SNsideheight}{\pgfpointdiff{\pgfpointanchor{SNsidebox}{south}}{\pgfpointanchor{SNsidebox}{north}}}\typeout{SNHEIGHT:24:22:\the\SNsideheight}}
\AddToShipoutPictureFG{\begin{tikzpicture}[remember picture,overlay]\ifcsname SNside\thepage\endcsname\csname SNside\thepage\endcsname\fi\end{tikzpicture}}